\documentclass[../DoAn.tex]{subfiles}
\begin{document}

Chương này trình bày quá trình thử nghiệm và đánh giá hệ thống đã xây dựng ở Chương~\ref{chapter:Implementation}: phương pháp và dữ liệu đánh giá, kết quả cho hai chức năng cốt lõi, khảo sát so sánh mô hình, phân tích độ nhạy của trọng số, thời gian phản hồi đo được và kiểm thử chức năng hệ thống.

\section{Đánh giá thực nghiệm}
\label{section:eval}
Do không có dữ liệu người dùng, hệ thống được đánh giá bằng nhiều tầng độ đo ngoại tuyến không vòng tròn, bổ sung bằng các kiểm chứng độc lập (nhịp độ thật, chuyển miền sang tập khác). Nguyên tắc xuyên suốt là độ đo cuối quyết định việc chấp nhận một thành phần, và mọi thay đổi phải vượt một cổng kiểm định không thụt lùi.

\subsection{Phương pháp và dữ liệu đánh giá}
\label{subsection:eval-method}
Tập đánh giá là toàn bộ \textbf{5.138 bài hát} của kho nhạc. Chức năng gợi ý theo màu được đánh giá trên \textbf{12 màu cơ bản} (theo Berlin \& Kay \cite{berlin1969basic}, cũng là 12 màu mà ICEAS \cite{jonauskaite2020universal} khảo sát) với mỗi truy vấn lấy $top\text{-}k=10$.

Vì không có tập nhạc Việt được con người gán nhãn màu--cảm xúc, \textbf{ground truth cho gợi ý theo màu không phải nhãn người} mà được tạo theo phương pháp hiệu chỉnh phân bố kho nhạc \cite{steck2018calibrated}: mỗi màu được ánh xạ thành một tọa độ Valence--Arousal, rồi quy về \textit{đích phân vị} tương ứng trong phân bố tích lũy (CDF) của kho nhạc. Độ đo chính là \textit{sai số nhắm trúng đích} (targeting error, TE) theo Phương trình~\eqref{eq:te}: khoảng cách Euclid trong không gian phân vị giữa đích phân vị của màu truy vấn và phân vị V-A của các bài được gợi ý, lấy trung bình trên các bài rồi trung bình trên 12 màu.
\begin{equation}\label{eq:te}
    \text{TE} = \frac{1}{|C|}\sum_{c \in C} \frac{1}{k}\sum_{s \in \text{top-}k(c)} \big\lVert \mathbf{q}_c - \mathbf{p}_s \big\rVert_2
\end{equation}
trong đó $C$ là tập 12 màu, $\mathbf{q}_c$ là đích phân vị của màu $c$ và $\mathbf{p}_s$ là phân vị V-A của bài $s$. Đây là một độ đo ngoại tuyến và \textit{tự tham chiếu} (đo trên chính nhãn V-A dùng để khớp); hạn chế này được nêu rõ ở cuối mục. Để tránh kết luận chỉ dựa vào một độ đo tự tham chiếu, TE được bổ sung bằng các kiểm chứng độc lập với nhãn (nhịp độ thật, chuẩn ICEAS, chuyển miền PMEmo) trình bày bên dưới. Các số liệu trong mục này được sinh bởi bộ công cụ đánh giá ngoại tuyến (\texttt{tools/color\_eval\_rigor.py}, \texttt{tools/color\_ablation.py}) và có thể chạy lại để tái lập.

\textbf{Phương pháp so sánh và lựa chọn mô hình.} Với mỗi vị trí có thể thay mô hình (trục biểu diễn âm thanh, embedding lời, các đầu dò cảm xúc, mô hình màu--cảm xúc), việc lựa chọn không dựa vào thứ hạng trên các bộ chuẩn quốc tế --- vì ``tốt trên giấy'' không bảo đảm tốt cho nhiệm vụ và dữ liệu cụ thể của hệ thống --- mà tuân theo năm nguyên tắc. Thứ nhất, (i) phán quyết bằng \textit{độ đo cuối} gắn với chính nhiệm vụ: NDCG trên ground truth tương đồng âm nhạc cho gợi ý bài tương tự, TE cùng các kiểm chứng độc lập cho gợi ý theo màu, và $R^2$ kiểm chứng chéo trên DEAM cho các đầu dò cảm xúc. Thứ hai, (ii) \textit{so sánh có kiểm soát}: chỉ thay đúng thành phần cần so, giữ nguyên mọi phần còn lại. Thứ ba, (iii) \textit{tinh chỉnh lại siêu tham số riêng} cho từng ứng viên trước khi so, vì kế thừa thiết lập của mô hình cũ có thể khiến một mô hình thực sự tốt hơn lại bị đánh giá là kém. Thứ tư, (iv) bảo đảm tính vững và chống hiện tượng ``tối ưu hóa chính độ đo'' bằng kiểm chứng chéo $k$-fold (tinh chỉnh trên fold huấn luyện, chấm trên fold tách riêng), khoảng tin cậy bootstrap và kiểm định ý nghĩa (Wilcoxon kèm hiệu chỉnh FDR), kèm \textit{tam giác hóa} bằng một tín hiệu độc lập (chuyển miền sang tập khác hoặc một bộ chấm độc lập). Cuối cùng, (v) chỉ chấp nhận thay thế khi ứng viên \textit{thắng hoặc hòa} trên độ đo cuối và vượt \textit{cổng không thụt lùi}, ưu tiên mô hình gọn/nhẹ hơn khi hòa. Kết quả áp dụng quy trình này cho toàn bộ các thành phần được trình bày ở mục~\ref{subsection:eval-bakeoff}.

\subsection{Đánh giá chức năng gợi ý theo màu}
\label{subsection:eval-color}
Bằng chứng \textit{chính} cho chức năng này là các kiểm chứng độc lập với nhãn (khớp chuẩn người ICEAS và nhịp độ thật, trình bày bên dưới); sai số nhắm đích (TE) --- vì tự tham chiếu --- chỉ đóng vai trò một \textit{điều kiện cần}. Cấu hình hiện tại đạt \textbf{TE $= 0.0242$} với khoảng tin cậy bootstrap 95\% là $[0.0219,\,0.0263]$. Bảng~\ref{table:phương pháp nền} và Hình~\ref{fig:te} so sánh với năm phương pháp nền. Hệ thống vượt cả bốn phương pháp nền thực (ngẫu nhiên, theo độ phổ biến, chỉ Valence, chỉ Arousal): kiểm định Wilcoxon kèm hiệu chỉnh tỉ lệ phát hiện sai Benjamini--Hochberg \cite{benjamini1995fdr} cho \textbf{4/4 bác bỏ giả thuyết không với $p$ hiệu chỉnh $= 0.00031$}. Hệ thống không vượt phương pháp nền ``lân cận V-A tham lam'' --- đây là một \textit{trần lý thuyết} (oracle) cố tình loại bỏ đa dạng nên không phải đối thủ công bằng; việc kết quả \textit{tiệm cận} trần này cho thấy cơ chế đa dạng hóa chỉ làm tăng TE rất ít.

\begin{longtable}{|p{6.6cm}|p{3.0cm}|p{3.2cm}|}
\caption{Sai số nhắm đích (TE) của hệ thống so với các phương pháp nền}\label{table:phương pháp nền}\\
\hline
\textbf{Phương pháp} & \textbf{TE} & \textbf{Ghi chú} \\
\hline
\endfirsthead
\hline
\textbf{Phương pháp} & \textbf{TE} & \textbf{Ghi chú} \\
\hline
\endhead
Ngẫu nhiên & 0.561 & phương pháp nền \\
\hline
Theo độ phổ biến & 0.513 & phương pháp nền \\
\hline
Chỉ Valence & 0.438 & loại bỏ một trục \\
\hline
Chỉ Arousal & 0.292 & loại bỏ một trục \\
\hline
Lân cận V-A tham lam (oracle) & 0.021 & trần lý thuyết, không đa dạng \\
\hline
\textbf{Brightify (V-A RBF)} & \textbf{0.0242} & CI95 $[0.0219, 0.0263]$ \\
\hline
\end{longtable}

\begin{figure}[tbp]
\centering
\includegraphics[width=0.78\textwidth]{Hinhve/fig_te_baselines.png}
\caption{So sánh TE của Brightify với các phương pháp nền (thấp hơn là tốt hơn); Brightify tiệm cận trần lý thuyết (oracle)}
\label{fig:te}
\end{figure}

Các kiểm chứng độc lập đều theo đúng hướng kỳ vọng. Valence của màu khớp chuẩn ICEAS với hệ số tương quan \textbf{$0.969$} (khoảng tin cậy Fisher-z $[0.889, 0.991]$); do chỉ có $n=12$ màu nên khoảng tin cậy rộng, và kiểm chứng giữ-một-ra (LOO-CV) cho $r \approx 0.87$. Công thức Arousal của Whiteford được kiểm bằng nhịp độ thật: $\rho(\text{độ sáng}, \text{BPM}) = 0.41$ và $\rho(\text{độ bão hòa}, \text{BPM}) = 0.66$. Hình~\ref{fig:colors} cho thấy 12 màu được ánh xạ vào các vùng cảm xúc khác nhau trên mặt phẳng V-A. Bảng~\ref{table:ablation} và Hình~\ref{fig:ablation} cho thấy đóng góp của hai cải tiến chính: so khớp trong không gian phân vị nâng \textbf{độ tách biệt giữa các màu từ $0.13$ lên $0.31$}, còn việc trừ trung bình khi đo nhất quán âm thanh nâng \textbf{độ nhất quán từ khoảng $0.02$ lên $0.23$}. Với hành trình hai màu, kiểm định Kolmogorov--Smirnov so với phân bố đều cho $\text{KS} = 0.21$ (đạt) và độ chênh nhịp độ giữa các bước chỉ ở mức $\approx 5.1$ BPM (chuỗi chuyển mượt).

\begin{figure}[tbp]
\centering
\begin{subfigure}{0.5\textwidth}
\centering
\includegraphics[width=\textwidth]{Hinhve/fig_colors_va.png}
\caption{12 màu trên mặt phẳng V-A}
\label{fig:colors}
\end{subfigure}\hfill
\begin{subfigure}{0.48\textwidth}
\centering
\includegraphics[width=\textwidth]{Hinhve/fig_ablation.png}
\caption{Đóng góp tích lũy (ablation)}
\label{fig:ablation}
\end{subfigure}
\caption{Trái: 12 màu kéo về các vùng cảm xúc khác nhau; Phải: mỗi cải tiến nâng độ tách biệt và độ nhất quán}
\label{fig:color-eval}
\end{figure}

\begin{longtable}{|p{5.4cm}|c|c|c|}
\caption{Ablation: mỗi cải tiến đóng góp vào tách biệt và nhất quán (12 màu, $top\text{-}k=10$)}\label{table:ablation}\\
\hline
\textbf{Cấu hình (cộng dồn)} & \textbf{Tách biệt} & \textbf{Nhất quán} & \textbf{TE} \\
\hline
\endfirsthead
\hline
\textbf{Cấu hình (cộng dồn)} & \textbf{Tách biệt} & \textbf{Nhất quán} & \textbf{TE} \\
\hline
\endhead
So khớp phân vị (cơ sở) & 0.129 & 0.015 & 0.0217 \\
\hline
+ Đích theo CDF & 0.225 & 0.026 & 0.0229 \\
\hline
+ Nhất quán âm thanh (trừ trung bình) & 0.224 & 0.113 & 0.0213 \\
\hline
+ Arousal Whiteford & 0.314 & 0.228 & 0.0242 \\
\hline
\end{longtable}

Các kết quả này có một số hạn chế cần nêu thẳng: (i) TE là độ đo \textit{tự tham chiếu} (offline), chưa có nghiên cứu nghe thực tế của người dùng; (ii) ánh xạ màu--cảm xúc là \textit{ngoại suy}: liên kết màu--cảm xúc tuy có khuôn mẫu phổ quát nhưng vẫn mang tính đặc thù văn hóa \cite{jonauskaite2019ml}, trong khi Việt Nam không nằm trong 30 quốc gia của khảo sát ICEAS; trục Valence--Arousal được ghi nhận ổn định xuyên văn hóa \cite{lee2025globalmood} song chưa có chuẩn màu--cảm xúc bản địa cho Việt Nam; (iii) màu ấm và bão hòa (nâu, hồng) được đọc là ``năng lượng cao'' --- điều này \textit{trung thành} với quan hệ màu--âm nhạc của Whiteford dù có thể nghịch trực giác.

\subsection{Đánh giá chức năng gợi ý bài tương tự}
\label{subsection:eval-similar}
Chức năng này được đánh giá trên các độ đo nội tại (80 bài hạt giống) và trên độ lợi NDCG \cite{steck2018calibrated} đối chiếu với các danh sách phát biên tập, theo tinh thần dùng phương pháp nền mạnh trong đánh giá hệ gợi ý \cite{dacrema2019progress}. Độ nhất quán tâm trạng (MoodCoherence) đo mức các bài kết quả cụm chặt trong không gian V-A, định nghĩa theo Phương trình~\eqref{eq:moodcoh}; độ nhất quán nhịp độ (TempoCoherence) định nghĩa tương tự trên nhịp độ.
\begin{equation}\label{eq:moodcoh}
    \text{MoodCoherence} = 1 - \frac{2}{k(k-1)}\sum_{i<j} \big\lVert \mathbf{p}_i - \mathbf{p}_j \big\rVert_2
\end{equation}
Về độ đo nội tại (80 bài hạt giống, $top\text{-}k=10$), MoodCoherence đạt $\approx 0.94$ và TempoCoherence đạt $\approx 0.84$ --- các bài kết quả cụm chặt cả về cảm xúc lẫn nhịp độ.

\textbf{Ground truth tương đồng âm nhạc (độ đo chính).} Để đánh giá đúng mức ``nghe giống nhau'', đồ án xây dựng một tập ground truth \textit{tương đồng âm nhạc} gồm 52 bài hạt giống, mỗi cặp được chấm theo thang liên quan $0$--$3$. Với mỗi bài hạt giống (lấy phân tầng theo cảm xúc để phủ đều các vùng Valence--Arousal), tập ứng viên được gộp theo lối TREC: hợp của nhóm bài hệ thống xếp hạng cao nhất và một nhóm bài lấy ngẫu nhiên làm đối chứng, tổng cộng $1.048$ cặp. Bộ chấm được hướng dẫn chỉ xét \textit{chất nhạc} (thể loại, nhịp--tiết tấu, mức năng lượng, độ ``nhảy'', giọng trưởng/thứ, không khí tổng thể), coi lời là gợi ý phụ và \textit{bỏ qua nghệ sĩ lẫn độ phổ biến}; nhờ vậy hạn chế được thiên lệch theo cùng nghệ sĩ, và phía hệ thống cũng đa dạng hóa nghệ sĩ kèm lọc bản cover nên tỉ lệ kết quả cùng nghệ sĩ ở mức thấp ($\approx 0.05$). Một lo ngại với ground truth do mô hình sinh là nó có thể là tạo tác của \textit{một} bộ chấm duy nhất; vì vậy toàn bộ 1.048 cặp được chấm lại độc lập bằng \textit{hai mô hình của hai nhà cung cấp khác nhau} (qwen3:8b cục bộ và GPT-4o-mini). Hai bộ chấm đồng thuận gần như tuyệt đối (Bảng~\ref{table:multijudge}): Spearman $\rho = 0.946$, đồng thuận nhị phân $0.99$ và Cohen $\kappa = 0.958$ --- cho thấy ground truth không phải hiện tượng riêng của một bộ chấm duy nhất. Trên tập này, hệ thống đạt NDCG@10 kiểm-chứng-chéo $\approx 0.73$ (chi tiết khảo sát backbone ở mục~\ref{subsection:eval-bakeoff}). Việc dùng mô hình ngôn ngữ làm bộ chấm là một lựa chọn \textit{thay thế} cho gán nhãn người trong điều kiện không tồn tại tập tương đồng nhạc Việt do người chấm; cách làm này theo tiền lệ gần đây trong đánh giá tương đồng âm nhạc và đã được giảm rủi ro bằng đồng thuận hai nhà cung cấp ở trên, song một nghiên cứu nghe với người thật vẫn là kiểm chứng cần thiết (Chương~\ref{chapter:conclusion}). Ngoài ra, dù tiếp cận theo nội dung giúp giảm phụ thuộc dữ liệu người dùng, hệ thống không hoàn toàn miễn nhiễm \textit{cold-start}: một bài hoặc nghệ sĩ mới chưa có đặc trưng ngoại tuyến sẽ chưa được gợi ý cho tới khi chạy lại pha trích đặc trưng.

\textbf{NDCG trên danh sách phát biên tập (độ đo phụ, đọc kèm phương pháp nền).} Trên 872 truy vấn từ danh sách phát biên tập, đo bằng độ lợi tích lũy chiết khấu chuẩn hóa NDCG \cite{jarvelin2002ndcg} (Phương trình~\eqref{eq:ndcg}), hệ sản phẩm đạt NDCG@10 $\approx 0.066$ (tái lập bằng \texttt{tools/eval\_editorial\_ndcg.py}). Tuy nhiên độ đo này \textit{không phù hợp} làm thước đo chính cho gợi ý bài tương tự: một phương pháp nền \textit{theo độ phổ biến} (chọn 10 bài theo tần suất nghệ sĩ) đạt $0.091$ --- \textit{cao hơn} hệ sản phẩm --- trong khi phương pháp nền ngẫu nhiên chỉ đạt $0.036$ (Bảng~\ref{table:editorial}). Đây không phải do hệ kém: danh sách phát biên tập \textit{thưởng} cho độ phổ biến và đồng-xuất-hiện, đồng thời \textit{phạt} chính cơ chế đa dạng nghệ sĩ (MMR) và lọc bản cover mà hệ cố ý áp dụng --- đúng như cảnh báo về việc dùng độ đo đại diện yếu trong đánh giá hệ gợi ý \cite{dacrema2019progress}. Do đó đồ án dùng ground truth tương đồng âm nhạc cùng các độ đo nội tại làm thước đo chính, còn NDCG biên tập chỉ để tham khảo xu hướng.
\begin{equation}\label{eq:ndcg}
    \text{NDCG@}k = \frac{\text{DCG@}k}{\text{IDCG@}k}, \qquad
    \text{DCG@}k = \sum_{i=1}^{k} \frac{rel_i}{\log_2(i+1)}
\end{equation}

\begin{longtable}{|p{8.0cm}|c|}
\caption{Đồng thuận hai bộ chấm độc lập trên ground truth tương đồng âm nhạc (1.048 cặp)}\label{table:multijudge}\\
\hline
\textbf{Chỉ số đồng thuận qwen3:8b $\leftrightarrow$ GPT-4o-mini} & \textbf{Giá trị} \\
\hline
\endfirsthead
\hline
\textbf{Chỉ số đồng thuận qwen3:8b $\leftrightarrow$ GPT-4o-mini} & \textbf{Giá trị} \\
\hline
\endhead
Spearman $\rho$ (thang liên quan $0$--$3$) & 0.946 \\
\hline
Đồng thuận nhị phân (liên quan $\geq 2$) & 0.99 \\
\hline
Cohen $\kappa$ & 0.958 \\
\hline
\end{longtable}

\begin{longtable}{|p{6.0cm}|c|c|}
\caption{NDCG@10 trên danh sách phát biên tập kèm phương pháp nền (đọc kèm cảnh báo Dacrema)}\label{table:editorial}\\
\hline
\textbf{Phương pháp} & \textbf{NDCG@10} & \textbf{$\times$ ngẫu nhiên} \\
\hline
\endfirsthead
\hline
\textbf{Phương pháp} & \textbf{NDCG@10} & \textbf{$\times$ ngẫu nhiên} \\
\hline
\endhead
Theo độ phổ biến (phương pháp nền) & 0.091 & 2.5$\times$ \\
\hline
Sản phẩm (Brightify) & 0.066 & 1.8$\times$ \\
\hline
Ngẫu nhiên (phương pháp nền) & 0.036 & 1.0$\times$ \\
\hline
\end{longtable}

Khả năng chuyển miền của các đầu dò cảm xúc được kiểm trên tập PMEmo \cite{zhang2018pmemo} (Hình~\ref{fig:pmemo}), đạt tương quan $0.69$ cho Valence và $0.65$ cho Arousal, cho thấy đầu dò nắm bắt cảm xúc thật chứ không phải hiện tượng riêng của tập huấn luyện; một kiểm chứng arousal độc lập bằng đặc trưng âm thanh khác (CLAP \cite{wu2023clap}) cho tương quan $0.46$.

\begin{figure}[tbp]
\centering
\includegraphics[width=0.5\textwidth]{Hinhve/fig_pmemo.png}
\caption{Chuyển miền các đầu dò cảm xúc sang tập PMEmo độc lập (tương quan Spearman, kèm khoảng tin cậy)}
\label{fig:pmemo}
\end{figure}

\subsection{Khảo sát so sánh mô hình}
\label{subsection:eval-bakeoff}
Để kiểm tra mỗi thành phần học máy có phải là lựa chọn tốt theo độ đo cuối hay không (trong phạm vi các ứng viên đã khảo sát), đồ án tiến hành một khảo sát so sánh trên \textbf{năm vị trí có thể thay mô hình}, với tổng cộng hơn mười mô hình tiên tiến được tải về và chạy thực tế; mỗi ứng viên được tinh chỉnh siêu tham số riêng trước khi so và được phán quyết bằng độ đo cuối theo phương pháp ở mục~\ref{subsection:eval-method}. Bảng~\ref{table:bakeoff} tóm tắt toàn bộ khảo sát, còn Bảng~\ref{table:backbone} trình bày sâu một head-to-head công bằng tiêu biểu (MuQ so với MERT). Điểm đáng chú ý xuyên suốt là phần lớn các mô hình ``mạnh trên giấy'' theo bộ chuẩn quốc tế lại \textit{không} vượt được độ đo cuối của hệ thống: MuQ-MuLan tối ưu cho việc ghép văn bản--nhạc nên cho sai số màu và độ nhất quán kém hơn MuQ; multilingual-e5-large \cite{wang2024e5} thắng các mô hình embedding lời còn lại (kể cả mô hình tiếng Việt tiền nhiệm) về độ nhất quán nội tại ở mọi trọng số; PhoBERT-large chỉ ngang ViSoBERT (ViSoBERT gọn hơn và hợp văn phong lời nhạc hơn); mDeBERTa-v3 kém XLM-R ở vai trò đầu dò liên ngữ; và với trục Arousal của màu, mô hình màu--nhạc của Whiteford được chọn thay cho Valdez--Mehrabian và bản hồi quy theo chuẩn ICEAS vì hai bản sau khớp arousal \textit{màu khi nhìn} (coi màu tối là kích thích) --- cấu trúc sai cho một hệ gợi ý nhạc. Kết quả này củng cố nguyên tắc đánh giá bằng độ đo cuối thay vì thứ hạng trên bộ chuẩn. Ngoài các thành phần phục vụ, hệ thống còn dùng CLAP làm một kiểm chứng arousal độc lập (mục~\ref{subsection:eval-similar}) và đã thử nhưng \textit{loại} bộ đặc trưng Essentia do đặc trưng suy biến ở tần số lấy mẫu 44.1\,kHz. Chi tiết phân tích khảo sát này được trình bày ở Chương~\ref{chapter:SolutionAndContribution}.

\textbf{Lựa chọn trục biểu diễn âm thanh (MuQ so với MERT).} Việc chọn MuQ thay cho MERT được kiểm bằng một so sánh \textit{công bằng}: chuyển backbone thật và \textit{tối ưu lại trọng số tổ hợp riêng cho từng backbone} bằng kiểm chứng chéo 5-fold trên ground truth tương đồng âm nhạc (Bảng~\ref{table:backbone}). Trên các độ đo trực tiếp, hai backbone gần như \textit{ngang nhau} (NDCG kiểm-chứng-chéo $0.735$ so với $0.734$; NDCG chỉ-âm-thanh $0.748$ so với $0.742$ --- chênh lệch không có ý nghĩa thống kê), và sai số màu TE coi như hòa. Khác biệt có ý nghĩa nằm ở \textit{phép chuyển giao sang một bộ chấm độc lập} (GPT-4o-mini): ở đó MuQ vượt MERT $0.806$ so với $0.771$ (khoảng tin cậy bootstrap cặp $[+0.014, +0.056]$, có ý nghĩa thống kê). Cùng với việc MuQ thắng ở đầu dò Arousal trên DEAM ($0.69$ so với $0.647$, mục~\ref{subsection:4.2.1}), kết quả cho thấy MuQ \textit{ngang hoặc hơn} MERT ở mọi mặt và việc chọn MuQ là hợp lý. Lưu ý phương pháp: trọng số tổ hợp tối ưu theo kiểm chứng chéo hội tụ về gần như chỉ-âm-thanh; trọng số sản phẩm $0.76/0.16/0.08$ đánh đổi rất ít NDCG để tăng độ nhất quán cảm xúc.

\begin{longtable}{|p{7.6cm}|c|c|}
\caption{So sánh công bằng MuQ với MERT (mỗi backbone tối ưu trọng số riêng)}\label{table:backbone}\\
\hline
\textbf{Phép đo} & \textbf{MuQ} & \textbf{MERT} \\
\hline
\endfirsthead
\hline
\textbf{Phép đo} & \textbf{MuQ} & \textbf{MERT} \\
\hline
\endhead
NDCG@10 kiểm-chứng-chéo (GT tương đồng âm nhạc) & 0.735 & 0.734 \\
\hline
NDCG@10 chỉ-âm-thanh (cô lập backbone) & 0.748 & 0.742 \\
\hline
NDCG@10 chuyển sang bộ chấm độc lập (GPT) & \textbf{0.806} & 0.771 \\
\hline
Tương quan Arousal kiểm-chứng-chéo trên DEAM & \textbf{0.69} & 0.647 \\
\hline
\end{longtable}

\begin{longtable}{|p{2.2cm}|p{2.5cm}|p{3.3cm}|p{4.5cm}|}
\caption{Khảo sát so sánh mô hình trên năm vị trí (mô hình được chọn in đậm)}\label{table:bakeoff}\\
\hline
\textbf{Vị trí} & \textbf{Được chọn} & \textbf{Ứng viên đã so sánh} & \textbf{Phán quyết trên độ đo cuối} \\
\hline
\endfirsthead
\hline
\textbf{Vị trí} & \textbf{Được chọn} & \textbf{Ứng viên đã so sánh} & \textbf{Phán quyết trên độ đo cuối} \\
\hline
\endhead
Biểu diễn âm thanh & \textbf{MuQ} & MERT, MuQ-MuLan & MuQ $\geq$ MERT (Bảng~\ref{table:backbone}); MuLan kém sai số màu và nhất quán vì tối ưu ghép văn bản--nhạc \\
\hline
Embedding lời & \textbf{multilingual\hspace{0pt}-e5-large} & VN-SBERT (tiền nhiệm), BGE-M3, PhoBERT-large, mDeBERTa & Thắng độ nhất quán nội tại ở 200 hạt giống, mọi trọng số \\
\hline
Cảm xúc lời (VN) & \textbf{ViSoBERT} & PhoBERT-large, wonrax-sentiment, ViDeBERTa & PhoBERT-large chỉ ngang (sau khi bật tách từ); ViSoBERT gọn và hợp văn phong lời hơn \\
\hline
Cảm xúc liên ngữ & \textbf{XLM-R} & mDeBERTa-v3 & mDeBERTa + gộp trung bình vẫn kém CCC \\
\hline
Màu $\rightarrow$ Arousal & \textbf{Whiteford 2018} & Valdez--Mehrabian, hồi quy ICEAS & Whiteford đúng cấu trúc màu--nhạc (kiểm bằng BPM thật); hai bản kia khớp arousal màu khi nhìn nên màu tối hóa ``năng lượng cao'' sai \\
\hline
\end{longtable}

\subsection{Phân tích độ nhạy của trọng số}
\label{subsection:eval-sensitivity}
Một phần trọng số trong hệ thống không được khớp trực tiếp từ dữ liệu người gán mà được \textit{đặt theo nguyên tắc} (trọng số nhịp độ trong Arousal, băng thông $\sigma$ của RBF màu) hoặc \textit{đánh đổi có chủ đích} giữa hai mục tiêu (trọng số V-A trong tổ hợp gợi ý bài tương tự). Để chứng minh các giá trị đang phục vụ không phải lựa chọn tùy tiện, đồ án thực hiện một phân tích độ nhạy: quét quanh mỗi giá trị và quan sát độ đo cuối thay đổi ra sao. Phân tích này \textit{không} thay đổi gì trong hệ đang chạy --- nó chỉ đo trên bộ nhớ và có thể tái lập bằng \texttt{tools/tune\_muq\_arousal.py} và \texttt{tools/sensitivity\_analysis.py}.

\textbf{Trọng số nhịp độ (Arousal).} Đây là một \textit{đánh đổi}, không phải điểm phẳng: tăng trọng số nhịp độ giúp Arousal bám nhịp độ thật tốt hơn nhưng làm giảm độ chính xác kiểm-chứng-chéo trên DEAM (Hình~\ref{fig:sens-tempo}). Giá trị phục vụ $0.35$ nằm trong vùng hợp lệ: nó đạt $\rho(\text{Arousal},\text{BPM}) = 0.47$ (vượt xa mục tiêu $0.20$) trong khi vẫn giữ DEAM-CV $= 0.69$, cao hơn cổng không-thụt-lùi $0.647$. Tăng lên $0.45$ sẽ phá cổng này (DEAM-CV tụt còn $0.61$), nên cửa sổ hợp lệ là khoảng $[0.25,\,0.40]$ và $0.35$ được chọn để bám nhịp độ mạnh nhất trong cửa sổ đó.

\begin{figure}[tbp]
\centering
\includegraphics[width=0.62\textwidth]{Hinhve/fig_sens_tempo.png}
\caption{Độ nhạy của trọng số nhịp độ: vùng xanh là cửa sổ hợp lệ (DEAM-CV $\geq 0.647$ và $\rho(A,\text{BPM}) \geq 0.20$); giá trị phục vụ $0.35$ nằm trong vùng này}
\label{fig:sens-tempo}
\end{figure}

\textbf{Băng thông $\sigma$ của RBF màu.} Ngược lại, kết quả gợi ý theo màu \textit{rất bền} với lựa chọn $\sigma$: khi quét $\sigma_V$ trong $[0.12,\,0.28]$ và $\sigma_A$ trong $[0.08,\,0.20]$, sai số nhắm đích (TE) chỉ dao động chưa tới $0.003$ (Hình~\ref{fig:sens-sigma}). Điều này cho thấy giá trị phục vụ ($\sigma_V = 0.20$, $\sigma_A = 0.14$) --- được đặt theo tỉ lệ độ tin cậy của hai trục (arousal đáng tin hơn nên $\sigma$ hẹp hơn) --- không phải tham số mong manh; kết luận không phụ thuộc vào con số chính xác.

\begin{figure}[tbp]
\centering
\includegraphics[width=0.62\textwidth]{Hinhve/fig_sens_sigma.png}
\caption{Độ nhạy của băng thông $\sigma$ (RBF màu): TE gần như phẳng trên toàn lưới quét, xác nhận hệ bền với lựa chọn $\sigma$}
\label{fig:sens-sigma}
\end{figure}

\textbf{Trọng số V-A trong gợi ý bài tương tự.} Đây là một đánh đổi được công bố thẳng. Khi tăng trọng số V-A từ $0$, NDCG@10 (trên ground truth tương đồng âm nhạc) đạt cao nhất ở mức gần \textit{chỉ-âm-thanh} rồi giảm dần, trong khi độ nhất quán cảm xúc (MoodCoherence) tăng mạnh (Hình~\ref{fig:sens-fusion}, Bảng~\ref{table:sens-fusion}). Giá trị phục vụ $0.16$ đúng là điểm \textit{khuỷu}: đi từ $0$ lên $0.16$, MoodCoherence nhảy từ $0.83$ lên $0.95$ ($+0.12$) trong khi NDCG chỉ giảm $\approx 0.003$; vượt quá $0.16$ thì NDCG tiếp tục giảm còn MoodCoherence gần như bão hòa. Nói cách khác, $0.16$ thu gần như toàn bộ phần lợi về nhất quán cảm xúc với chi phí NDCG rất nhỏ. (Lưu ý: giá trị NDCG tuyệt đối ở đây thấp vì được tính bằng cách xếp hạng \textit{toàn bộ kho nhạc} rồi đối chiếu với tập đã chấm nhỏ, khác với NDCG kiểm-chứng-chéo trên pool ở mục~\ref{subsection:eval-bakeoff}; chỉ \textit{xu hướng tương đối} giữa các trọng số là điều cần đọc.)

\begin{longtable}{|c|c|c|c|}
\caption{Độ nhạy của trọng số V-A trong gợi ý bài tương tự (52 hạt giống ground truth tương đồng âm nhạc)}\label{table:sens-fusion}\\
\hline
\textbf{Trọng số V-A} & \textbf{Trọng số âm thanh} & \textbf{NDCG@10 (xu hướng)} & \textbf{MoodCoherence} \\
\hline
\endfirsthead
\hline
\textbf{Trọng số V-A} & \textbf{Trọng số âm thanh} & \textbf{NDCG@10 (xu hướng)} & \textbf{MoodCoherence} \\
\hline
\endhead
0.00 & 0.92 & 0.0414 & 0.833 \\
\hline
0.08 & 0.84 & 0.0362 & 0.924 \\
\hline
\textbf{0.16 (phục vụ)} & \textbf{0.76} & \textbf{0.0386} & \textbf{0.948} \\
\hline
0.24 & 0.68 & 0.0335 & 0.957 \\
\hline
0.32 & 0.60 & 0.0299 & 0.964 \\
\hline
\end{longtable}

\begin{figure}[tbp]
\centering
\includegraphics[width=0.66\textwidth]{Hinhve/fig_sens_fusion.png}
\caption{Đánh đổi NDCG $\leftrightarrow$ MoodCoherence theo trọng số V-A: giá trị phục vụ $0.16$ là điểm khuỷu --- gần như toàn bộ lợi ích nhất quán cảm xúc đạt được với chi phí NDCG nhỏ}
\label{fig:sens-fusion}
\end{figure}

Tổng hợp lại, phân tích độ nhạy cho thấy: ở những trục đặt theo nguyên tắc ($\sigma$), kết quả bền với lựa chọn chính xác; còn ở những trục mang tính đánh đổi (nhịp độ, V-A), giá trị phục vụ nằm đúng vùng/điểm cân bằng đã nêu lý do. Không có trọng số nào được chọn cảm tính mà không kiểm chứng vùng lân cận.

\subsection{Bàn luận và mối đe dọa tính hợp lệ}
\label{subsection:eval-discussion}
Tổng hợp lại, các kết quả trên cho thấy hệ thống đạt mục tiêu kép: khớp cảm xúc (TE thấp, vượt mọi phương pháp nền thực) và nghe đồng nhất (độ nhất quán tăng mạnh). Tuy nhiên, để diễn giải đúng mức ý nghĩa của chúng, cần nêu rõ các mối đe dọa tính hợp lệ theo ba nhóm thường dùng trong đánh giá hệ thống.

\textit{Hợp lệ về cấu trúc} (construct validity) --- liệu độ đo có đo đúng thứ cần đo. Sai số nhắm đích TE đo trên chính nhãn V-A được dùng để truy hồi, nên về bản chất là \textit{tự tham chiếu}: nó kiểm tra hệ thống có truy hồi đúng theo định nghĩa của chính nó hay không, chứ chưa khẳng định nhãn V-A là ``đúng'' theo cảm nhận người. Đây là lý do đồ án bổ sung các kiểm chứng \textit{độc lập với nhãn}: tương quan với nhịp độ thật (cho trục Arousal của màu), khớp chuẩn ICEAS (cho trục Valence của màu), và chuyển miền sang PMEmo (cho các đầu dò cảm xúc). Sự nhất quán giữa các kiểm chứng độc lập này làm tăng đáng kể độ tin cậy rằng nhãn V-A nắm bắt cảm xúc thật.

\textit{Nguyên tắc dùng dữ liệu nước ngoài.} Một câu hỏi tự nhiên là vì sao có thể dùng các tập do con người gán nhãn ở nước ngoài (DEAM, EmoBank, PMEmo, chuẩn màu ICEAS) cho một hệ thống nhạc Việt. Nguyên tắc xuyên suốt là dữ liệu nước ngoài chỉ được dùng ở những trục mà tín hiệu \emph{không phụ thuộc ngôn ngữ}. Arousal được suy từ năng lượng âm thanh (nhịp độ, độ to, đặc trưng MuQ) --- một bài nhanh và to mang năng lượng cao bất kể hát bằng tiếng nào --- nên đầu dò huấn luyện trên DEAM (nhạc phương Tây) chuyển sang nhạc Việt là hợp lệ; tương tự, ánh xạ màu dựa trên tri giác màu của chuẩn ICEAS. Ngược lại, Valence \emph{phụ thuộc nghĩa của lời} nên \emph{không} được lấy chủ yếu từ dữ liệu nước ngoài: trong tổ hợp EWE, các tín hiệu lời chiếm $\approx 91\%$ trọng số, trong đó $\approx 57\%$ đến từ tài nguyên \emph{tiếng Việt bản địa} (từ điển NRC-VAD-VN cùng VnEmoLex, và đầu dò ViSoBERT trên dữ liệu cảm xúc tiếng Việt) và $\approx 34\%$ từ EmoBank tiếng Anh được \emph{chuyển ngữ} qua mô hình đa ngữ XLM-R; đầu dò âm thanh chỉ đóng góp $\approx 9\%$. Cách phân vai theo phương thức (âm thanh/tri giác dùng nguồn quốc tế, ngữ nghĩa lời dùng nguồn bản địa) chính là điều được kiểm chứng bằng phép đo nêu ở mục hợp lệ ngoại tại: nếu cố suy Valence từ âm thanh nước ngoài thì kết quả rất yếu ($R^2 \approx 0.06$).

\textit{Vì sao không tự gán nhãn cảm xúc cho kho nhạc Việt.} Việc không xây một bộ nhãn cảm xúc do chính tác giả gán là một lựa chọn có chủ đích về phương pháp, không phải do thiếu dữ liệu. Thứ nhất, một bộ nhãn do \emph{một người gán duy nhất} không có độ tin cậy liên-người-gán (inter-rater reliability) nên về mặt thống kê không cấu thành một ground truth đáng tin --- nó chỉ phản ánh phán đoán của một cá nhân. Thứ hai, vì người gán cũng chính là người xây dựng hệ thống, nhãn tự gán có nguy cơ \emph{thiên lệch xác nhận} (confirmation bias) khiến việc ``kiểm chứng'' trở thành lập luận vòng tròn. Thay vì tạo ra một chuẩn so yếu và dễ thiên lệch như vậy, đồ án kiểm chứng nhãn bằng sự \emph{hội tụ với nhiều tập do con người gán nhãn đã công bố và bình duyệt} (DEAM, PMEmo cho âm thanh; ICEAS cho màu --- mỗi tập là phán đoán tổng hợp của hàng chục đến hàng trăm người gán độc lập) cùng một panel nhiều bộ chấm (mục~\ref{subsection:eval-similar}); riêng Valence còn được đối chiếu với một bộ tham chiếu GPT-4o-mini đọc trực tiếp lời tiếng Việt ($\rho = 0.63$, mục~\ref{subsection:4.2.1}). Sự nhất quán giữa các nguồn độc lập này là bằng chứng khách quan hơn so với một tập tự gán đơn lẻ. Việc xây dựng một bộ dữ liệu cảm xúc nhạc Việt có nhiều người gán và báo cáo độ tin liên-người-gán được để lại như một hướng phát triển (Chương~\ref{chapter:conclusion}).

\textit{Hợp lệ nội tại} (internal validity) --- liệu kết luận so sánh có vững. Việc so với năm phương pháp nền (trong đó có hai đường loại bỏ từng trục) kèm kiểm định Wilcoxon và hiệu chỉnh FDR bảo đảm cải thiện không phải do ngẫu nhiên; khoảng tin cậy bootstrap cho thấy độ ổn định của ước lượng. Riêng đường ``lân cận V-A tham lam'' là một trần lý thuyết (oracle) chứ không phải đối thủ công bằng, nên việc hệ thống tiệm cận nó (chứ không vượt) là kết quả mong đợi, cho thấy phần đa dạng hóa chỉ đánh đổi rất ít độ chính xác nhắm đích. Riêng với gợi ý bài tương tự, ground truth tương đồng âm nhạc còn được xác nhận bằng hai bộ chấm độc lập của hai nhà cung cấp ($\kappa = 0.958$, mục~\ref{subsection:eval-similar}), nên kết luận so sánh không phải là tạo tác của một bộ chấm duy nhất; hạn chế còn lại là cỡ mẫu hạt giống ($n = 52$).

\textit{Hợp lệ ngoại tại} (external validity) --- liệu kết quả có tổng quát hóa. Đây là hạn chế lớn nhất: chưa có nghiên cứu nghe của người dùng thật (theo Dacrema và cộng sự, tương quan giữa độ đo ngoại tuyến và mức hài lòng thực tế chỉ vào khoảng $r \approx 0.28$), và ánh xạ màu--cảm xúc là ngoại suy từ dữ liệu quốc tế (Việt Nam không nằm trong khảo sát ICEAS gốc). Một hạn chế cụ thể của trục Valence là \textit{chuyển giao yếu qua biên giới văn hóa}: trong phạm vi nhạc phương Tây, phép chuyển miền DEAM$\to$PMEmo khá tốt ($\rho = 0.69$ cho Valence), nhưng khi áp đầu dò âm thanh huấn luyện trên DEAM sang nhạc Việt thì Valence rất yếu ($R^2 \approx 0.06$). Đây chính là lý do Valence được lấy chủ yếu từ các tín hiệu lời (phần lớn là tài nguyên tiếng Việt bản địa, phần còn lại là EmoBank tiếng Anh chuyển ngữ qua XLM-R; chi tiết phân vai và trọng số nêu ở phần hợp lệ cấu trúc trên) thay vì từ âm thanh nước ngoài --- một lựa chọn đúng nhưng phụ thuộc chất lượng các tài nguyên đó và độ tốt của phép chuyển ngữ liên ngữ. Ngoài ra, do tương quan dương giữa Valence và Arousal trong kho nhạc, các góc phần tư lệch trục (``bình yên'': V cao--A thấp; ``căng thẳng'': V thấp--A cao) thưa dữ liệu, khiến màu ấm phải nhắm vào phần đuôi Arousal cao mỏng ($\approx 22\%$ kho nhạc) --- đây là thuộc tính của dữ liệu. Chỉ số NDCG@10 biên tập tuyệt đối thấp ($\approx 0.066$) cũng cần đọc thận trọng: nó là đại diện gần đúng cho mức liên quan và còn bị một phương pháp nền theo độ phổ biến vượt qua (mục~\ref{subsection:eval-similar}), nên chỉ phù hợp để tham khảo xu hướng chứ không khẳng định chất lượng tuyệt đối. Những hạn chế này được nêu lại và đề xuất hướng khắc phục ở Chương~\ref{chapter:conclusion}.

\section{Thời gian phản hồi và triển khai}
\label{section:deploy}
Hệ thống được triển khai thử nghiệm với pha ngoại tuyến chạy trước để sinh các đặc trưng và ma trận biểu diễn cho toàn bộ kho nhạc, sau đó động cơ gợi ý nạp các đặc trưng này vào bộ nhớ và phục vụ truy vấn qua giao diện lập trình ứng dụng cho ứng dụng trang đơn.

\subsection{Thời gian phản hồi đo được}
\label{subsec:latency}
Vì mọi tính toán nặng đã thực hiện ở pha ngoại tuyến và không có mô hình nào được suy luận lúc phục vụ (nhãn cảm xúc đã đóng băng thành tra cứu), mỗi truy vấn chỉ gồm vài phép nhân ma trận trên ma trận $5138 \times 1024$. Đo trực tiếp trong tiến trình trên một máy Apple Silicon bằng \texttt{tools/bench\_latency.py} (1000 lần lặp, chưa tính thời gian mạng và tuần tự hóa JSON), kết quả như Bảng~\ref{table:latency} và Hình~\ref{fig:latency}: nạp động cơ lõi một lần chỉ mất vài giây (đo được khoảng $1$--$6$ giây tùy bộ nhớ đệm; khởi động đầy đủ kèm bộ mã hóa tìm kiếm tùy chọn sẽ lâu hơn), còn mỗi truy vấn lúc phục vụ ở mức \textbf{vài mili-giây} (gợi ý theo một màu $\approx 3.1$ ms, gợi ý bài tương tự $\approx 4.0$ ms); hành trình hai màu nặng hơn ($\approx 70$ ms) do bước sắp xếp chuỗi cảm xúc. Đây là độ trễ tính toán lõi; độ trễ đầu cuối nhìn từ trình duyệt sẽ cao hơn do mạng và phát âm thanh.

\begin{longtable}{|p{6.6cm}|c|c|}
\caption{Thời gian phản hồi tính toán lõi (đo trong tiến trình, không gồm mạng)}\label{table:latency}\\
\hline
\textbf{Thao tác} & \textbf{Trung vị} & \textbf{p95} \\
\hline
\endfirsthead
\hline
\textbf{Thao tác} & \textbf{Trung vị} & \textbf{p95} \\
\hline
\endhead
Nạp động cơ lõi (một lần, lúc khởi động) & $\approx 1$--$6$ s & --- \\
\hline
Gợi ý theo một màu (UC02) & $3.1$ ms & $4.4$ ms \\
\hline
Gợi ý bài tương tự (UC01) & $4.0$ ms & $4.5$ ms \\
\hline
Hành trình hai màu & $70$ ms & $73.6$ ms \\
\hline
\end{longtable}

\begin{figure}[tbp]
\centering
\includegraphics[width=0.68\textwidth]{Hinhve/fig_latency.png}
\caption{Thời gian phản hồi tính toán lõi theo thao tác (thang log); truy vấn đơn ở mức vài mili-giây}
\label{fig:latency}
\end{figure}

\section{Kiểm thử chức năng hệ thống}
\label{section:functest}
Ngoài đánh giá mô hình, hệ thống còn được kiểm thử như một sản phẩm phần mềm. Bảng~\ref{table:testcase} tóm tắt các ca kiểm thử chính. Một phần được tự động hóa bằng \texttt{pytest} (\texttt{test/test\_color\_reco.py} kiểm tính hợp lệ cấu hình, ánh xạ 12 màu về đúng góc phần tư cảm xúc, đa dạng nghệ sĩ, tính đơn điệu của hành trình, độ trải Arousal; \texttt{test/test\_endless\_radio.py} kiểm hợp đồng loại trừ để chế độ ``đài'' không lặp), một phần được kiểm theo hợp đồng API và thủ công. Hai tệp cốt lõi này gồm 27 ca kiểm thử tự động; khi chạy toàn bộ bộ kiểm thử của dự án bằng \texttt{pytest}, có \textbf{124 ca đạt}.

\begin{longtable}{|p{1.0cm}|p{4.4cm}|p{5.2cm}|p{2.6cm}|}
\caption{Các ca kiểm thử chức năng hệ thống}\label{table:testcase}\\
\hline
\textbf{Mã} & \textbf{Mục tiêu} & \textbf{Kết quả mong đợi} & \textbf{Trạng thái} \\
\hline
\endfirsthead
\hline
\textbf{Mã} & \textbf{Mục tiêu} & \textbf{Kết quả mong đợi} & \textbf{Trạng thái} \\
\hline
\endhead
TC01 & Chọn màu và nhận danh sách bài hát & Trả về $\geq 5$ bài, V-A hợp lệ, đa dạng nghệ sĩ & Đạt (tự động) \\
\hline
TC02 & Chọn bài và nhận bài tương tự & Trả về danh sách không chứa bài nguồn, đã loại cover & Đạt (tự động) \\
\hline
TC03 & Xử lý mã màu không hợp lệ & Từ chối với lỗi xác thực đầu vào, không lỗi máy chủ & Đạt (hợp đồng API) \\
\hline
TC04 & Xử lý bài không có lời & Trường lời rỗng, không gây lỗi & Đạt (hợp đồng API) \\
\hline
TC05 & Xử lý bài không có đặc trưng & Trường tương ứng rỗng; lùi về tín hiệu còn lại & Đạt (hợp đồng API) \\
\hline
TC06 & Kiểm tra thời gian phản hồi & Truy vấn đơn ở mức vài mili-giây & Đạt (mục~\ref{subsec:latency}) \\
\hline
TC07 & Kiểm tra giao diện trên trình duyệt & Hai giao diện dựng và hoạt động & Đạt một phần (smoke; đa trình duyệt: hướng hoàn thiện) \\
\hline
\end{longtable}

Các kết quả đánh giá định lượng và kiểm thử trên là cơ sở để rút ra các giải pháp và đóng góp nổi bật được phân tích ở chương tiếp theo.

\end{document}
